% =================================================================
%  CSE 815 -- Machine Learning
%  Module 2, Week 6, Part 1
%  Advice for Applying Machine Learning: Evaluation, Model Selection,
%  Bias and Variance, and Learning Curves
% =================================================================
\documentclass[aspectratio=169,10pt]{beamer}
\usepackage{cse815}

\title[CSE 815 -- Week 6, Part 1]{Advice for Applying Machine Learning}
\subtitle{Evaluating Models, Choosing Among Them, and Diagnosing Bias and Variance}
\author{Rokan Uddin Faruqui\\ \small Professor, Department of CSE}
\institute{\small University of Chittagong \\ \texttt{rufaruqui@cu.ac.bd}}
\date[Week 6 / Part 1]{Week 6 $\cdot$ Part 1}

\begin{document}

\begin{frame}[plain,noframenumbering]
  \titlepage
\end{frame}

% ---------------------------------------------------------------
\begin{frame}{From building models to improving them}
  \begin{columns}[T,onlytextwidth]
    \begin{column}{0.5\textwidth}
      \begin{block}{Weeks 1--5 gave us}
        \footnotesize
        Linear and logistic regression, regularisation, and neural networks trained with
        Adam. You can now build a model for almost any supervised task.
      \end{block}
      \footnotesize
      Building is the easy half. When the first model disappoints --- and it usually does ---
      there are a dozen things you could try, and some of them cost months. This week is
      about choosing \emph{which} one, from evidence rather than hunches.
    \end{column}
    \begin{column}{0.46\textwidth}
      \begin{keyidea}[Today]
        \footnotesize
        \begin{enumerate}
          \item Deciding what to try next; evaluating a model
          \item Model selection: training, validation and test sets
          \item Diagnosing bias and variance; the role of $\lambda$
          \item Baselines and learning curves
          \item What to try next, revisited; neural networks
        \end{enumerate}
      \end{keyidea}
      \footnotesize Part 2 covers the development process itself: error analysis, adding
      data, transfer learning, deployment, fairness, and metrics for skewed classes.
    \end{column}
  \end{columns}
\end{frame}

% ===============================================================
\section{Evaluating a Model}

\begin{frame}{Your model makes large errors. What next?}
  \begin{columns}[T,onlytextwidth]
    \begin{column}{0.5\textwidth}
      \begin{block}{Things you could try}
        \footnotesize
        \begin{enumerate}
          \item Get more training examples.
          \item Try a smaller set of features.
          \item Get additional features.
          \item Add polynomial features ($x_1^2$, $x_2^2$, $x_1x_2$, \dots).
          \item Decrease the regularisation parameter $\lambda$.
          \item Increase $\lambda$.
        \end{enumerate}
      \end{block}
      \footnotesize
      Each is reasonable. Each helps some models and is useless for others. Collecting more
      data alone can take a team months.
    \end{column}
    \begin{column}{0.46\textwidth}
      \begin{defnbox}[Diagnostic]
        \small A test you run to learn what is and is not working in a learning algorithm,
        so that you can decide how to improve it.
      \end{defnbox}
      \begin{keyidea}
        Diagnostics take time to run, but much less time than pursuing the wrong fix. The
        rest of this lecture builds the two most useful ones: the \textbf{bias/variance}
        diagnosis and the \textbf{learning curve}.
      \end{keyidea}
    \end{column}
  \end{columns}
\end{frame}

\begin{frame}{The running example: a synthetic regression task}
  \begin{columns}[T,onlytextwidth]
    \begin{column}{0.52\textwidth}
      \centering {\scriptsize 20 training points; $y=\sin 2x+0.5x$ plus noise, $\sigma=0.25$}\par
      \vspace{1mm}
      \begin{tikzpicture}
        \begin{axis}[width=7.2cm,height=4.9cm,xmin=0.15,xmax=3.05,ymin=-0.8,ymax=2.2,
          xlabel={$x$},ylabel={$y$},tick label style={font=\scriptsize},
          label style={font=\scriptsize},
          legend style={font=\tiny,at={(0.5,1.03)},anchor=south,legend columns=4,draw=none},
          legend cell align=left]
          \addplot[only marks,mark=*,mark size=1.6pt,cublue] coordinates {(2.829,1.344)(1.534,0.714)(2.929,0.979)(0.243,0.702)(1.822,0.191)(1.129,1.245)(2.406,0.211)(0.524,1.321)(2.615,0.109)(1.632,1.037)(2.707,0.501)(1.431,1.033)(1.291,1.388)(2.367,0.349)(2.952,1.372)(1.109,1.395)(2.907,0.996)(2.787,0.822)(0.533,0.893)(1.827,0.727)};
          \addlegendentry{training data}
          \addplot[cuamber,thick] coordinates {(0.200,1.051)(0.225,1.048)(0.250,1.046)(0.275,1.043)(0.300,1.040)(0.325,1.037)(0.350,1.035)(0.375,1.032)(0.400,1.029)(0.425,1.026)(0.450,1.024)(0.475,1.021)(0.500,1.018)(0.525,1.015)(0.550,1.013)(0.575,1.010)(0.600,1.007)(0.625,1.004)(0.650,1.002)(0.675,0.999)(0.700,0.996)(0.725,0.993)(0.750,0.991)(0.775,0.988)(0.800,0.985)(0.825,0.982)(0.850,0.980)(0.875,0.977)(0.900,0.974)(0.925,0.971)(0.950,0.969)(0.975,0.966)(1.000,0.963)(1.025,0.960)(1.050,0.958)(1.075,0.955)(1.100,0.952)(1.125,0.949)(1.150,0.947)(1.175,0.944)(1.200,0.941)(1.225,0.938)(1.250,0.936)(1.275,0.933)(1.300,0.930)(1.325,0.927)(1.350,0.925)(1.375,0.922)(1.400,0.919)(1.425,0.916)(1.450,0.914)(1.475,0.911)(1.500,0.908)(1.525,0.905)(1.550,0.903)(1.575,0.900)(1.600,0.897)(1.625,0.894)(1.650,0.892)(1.675,0.889)(1.700,0.886)(1.725,0.883)(1.750,0.881)(1.775,0.878)(1.800,0.875)(1.825,0.872)(1.850,0.870)(1.875,0.867)(1.900,0.864)(1.925,0.861)(1.950,0.859)(1.975,0.856)(2.000,0.853)(2.025,0.850)(2.050,0.848)(2.075,0.845)(2.100,0.842)(2.125,0.839)(2.150,0.837)(2.175,0.834)(2.200,0.831)(2.225,0.828)(2.250,0.826)(2.275,0.823)(2.300,0.820)(2.325,0.817)(2.350,0.815)(2.375,0.812)(2.400,0.809)(2.425,0.806)(2.450,0.804)(2.475,0.801)(2.500,0.798)(2.525,0.795)(2.550,0.793)(2.575,0.790)(2.600,0.787)(2.625,0.784)(2.650,0.782)(2.675,0.779)(2.700,0.776)(2.725,0.773)(2.750,0.771)(2.775,0.768)(2.800,0.765)(2.825,0.762)(2.850,0.760)(2.875,0.757)(2.900,0.754)(2.925,0.751)(2.950,0.749)(2.975,0.746)(3.000,0.743)};
          \addlegendentry{degree 1}
          \addplot[cuteal,very thick] coordinates {(0.200,0.427)(0.225,0.511)(0.250,0.592)(0.275,0.668)(0.300,0.741)(0.325,0.809)(0.350,0.873)(0.375,0.934)(0.400,0.990)(0.425,1.043)(0.450,1.093)(0.475,1.138)(0.500,1.181)(0.525,1.219)(0.550,1.255)(0.575,1.287)(0.600,1.317)(0.625,1.343)(0.650,1.366)(0.675,1.386)(0.700,1.403)(0.725,1.418)(0.750,1.430)(0.775,1.439)(0.800,1.446)(0.825,1.450)(0.850,1.452)(0.875,1.452)(0.900,1.450)(0.925,1.445)(0.950,1.438)(0.975,1.430)(1.000,1.419)(1.025,1.407)(1.050,1.393)(1.075,1.377)(1.100,1.359)(1.125,1.341)(1.150,1.320)(1.175,1.299)(1.200,1.276)(1.225,1.252)(1.250,1.227)(1.275,1.201)(1.300,1.174)(1.325,1.146)(1.350,1.117)(1.375,1.087)(1.400,1.057)(1.425,1.027)(1.450,0.995)(1.475,0.964)(1.500,0.932)(1.525,0.900)(1.550,0.868)(1.575,0.836)(1.600,0.803)(1.625,0.771)(1.650,0.739)(1.675,0.707)(1.700,0.676)(1.725,0.645)(1.750,0.614)(1.775,0.584)(1.800,0.555)(1.825,0.526)(1.850,0.498)(1.875,0.471)(1.900,0.445)(1.925,0.420)(1.950,0.396)(1.975,0.373)(2.000,0.352)(2.025,0.331)(2.050,0.313)(2.075,0.295)(2.100,0.280)(2.125,0.266)(2.150,0.254)(2.175,0.243)(2.200,0.235)(2.225,0.228)(2.250,0.224)(2.275,0.221)(2.300,0.221)(2.325,0.223)(2.350,0.228)(2.375,0.235)(2.400,0.244)(2.425,0.256)(2.450,0.271)(2.475,0.288)(2.500,0.309)(2.525,0.332)(2.550,0.358)(2.575,0.388)(2.600,0.420)(2.625,0.456)(2.650,0.495)(2.675,0.538)(2.700,0.584)(2.725,0.633)(2.750,0.687)(2.775,0.743)(2.800,0.804)(2.825,0.869)(2.850,0.937)(2.875,1.010)(2.900,1.086)(2.925,1.167)(2.950,1.252)(2.975,1.342)(3.000,1.436)};
          \addlegendentry{degree 3}
          \addplot[curose,thick] coordinates {(0.200,-4.553)(0.225,-1.092)(0.250,1.334)(0.275,2.935)(0.300,3.888)(0.325,4.343)(0.350,4.425)(0.375,4.237)(0.400,3.864)(0.425,3.375)(0.450,2.824)(0.475,2.253)(0.500,1.695)(0.525,1.174)(0.550,0.706)(0.575,0.301)(0.600,-0.034)(0.625,-0.299)(0.650,-0.493)(0.675,-0.620)(0.700,-0.686)(0.725,-0.697)(0.750,-0.660)(0.775,-0.582)(0.800,-0.472)(0.825,-0.335)(0.850,-0.181)(0.875,-0.014)(0.900,0.158)(0.925,0.331)(0.950,0.500)(0.975,0.661)(1.000,0.811)(1.025,0.948)(1.050,1.068)(1.075,1.172)(1.100,1.258)(1.125,1.327)(1.150,1.378)(1.175,1.412)(1.200,1.430)(1.225,1.433)(1.250,1.423)(1.275,1.401)(1.300,1.369)(1.325,1.329)(1.350,1.281)(1.375,1.228)(1.400,1.172)(1.425,1.113)(1.450,1.054)(1.475,0.995)(1.500,0.938)(1.525,0.883)(1.550,0.831)(1.575,0.784)(1.600,0.740)(1.625,0.701)(1.650,0.667)(1.675,0.638)(1.700,0.613)(1.725,0.593)(1.750,0.577)(1.775,0.565)(1.800,0.556)(1.825,0.549)(1.850,0.545)(1.875,0.543)(1.900,0.541)(1.925,0.540)(1.950,0.539)(1.975,0.537)(2.000,0.534)(2.025,0.530)(2.050,0.523)(2.075,0.515)(2.100,0.503)(2.125,0.489)(2.150,0.473)(2.175,0.454)(2.200,0.432)(2.225,0.408)(2.250,0.382)(2.275,0.354)(2.300,0.326)(2.325,0.297)(2.350,0.269)(2.375,0.242)(2.400,0.218)(2.425,0.196)(2.450,0.179)(2.475,0.168)(2.500,0.163)(2.525,0.166)(2.550,0.178)(2.575,0.200)(2.600,0.233)(2.625,0.278)(2.650,0.335)(2.675,0.405)(2.700,0.486)(2.725,0.579)(2.750,0.680)(2.775,0.788)(2.800,0.897)(2.825,1.004)(2.850,1.100)(2.875,1.175)(2.900,1.218)(2.925,1.213)(2.950,1.142)(2.975,0.980)(3.000,0.701)};
          \addlegendentry{degree 10}
        \end{axis}
      \end{tikzpicture}
    \end{column}
    \begin{column}{0.44\textwidth}
      \footnotesize
      \begin{block}{Setting}
        50 points from a known curve plus Gaussian noise: 20 for training, 15 for validation,
        15 for testing. Models: polynomials of degree $d=1,\dots,10$. Cost
        $J=\frac{1}{2m}\sum_i(f(x^{(i)})-y^{(i)})^2$.
      \end{block}
      \begin{block}{What the picture shows}
        Degree 1 misses the curve. Degree 10 passes close to every training point and swings
        wildly between them --- off the chart near $x=0.3$. Degree 3 looks right.
      \end{block}
      \begin{pitfall}[Plots do not scale]
        \footnotesize With one feature we can look. With $64$ pixels or $40$ weather
        stations we cannot. We need a number that says ``degree 10 is bad''.
      \end{pitfall}
    \end{column}
  \end{columns}
\end{frame}

\begin{frame}{Evaluating on data the model has not seen}
  \begin{columns}[T,onlytextwidth]
    \begin{column}{0.5\textwidth}
      \begin{block}{Split the data}
        \footnotesize
        Train on one part (often 70--80\%); hold back the rest as a \textbf{test set}, used
        only to measure.
        \[ J_{\text{test}}=\frac{1}{2m_{\text{test}}}\sum_{i=1}^{m_{\text{test}}}
           \bigl(f(\vx_{\text{test}}^{(i)})-y_{\text{test}}^{(i)}\bigr)^2 \]
        Compute it \emph{without} the regularisation term.
      \end{block}
      \begin{block}{For classification}
        \footnotesize
        $J_{\text{test}}$ can be the logistic cost, but the more readable measure is the
        \textbf{misclassification error}: the fraction of test examples with
        $\yhat\ne y$.
      \end{block}
    \end{column}
    \begin{column}{0.46\textwidth}
      \footnotesize
      \begin{block}{Our example, degree 10}
        \renewcommand{\arraystretch}{1.15}
        \begin{tabular}{@{}lr@{}}
          $J_{\text{train}}$ & $0.0164$\\
          $J_{\text{test}}$  & $0.3362$\\
        \end{tabular}\par
        \vspace{1mm}
        A training cost this low says the curve threads the training points. The test cost,
        twenty times larger, says it has memorised them.
      \end{block}
      \begin{keyidea}
        $J_{\text{train}}$ measures fit. Only error on held-out data estimates how the model
        will do on new examples --- its \textbf{generalisation error}.
      \end{keyidea}
    \end{column}
  \end{columns}
\end{frame}

% ===============================================================
\section{Model Selection and Cross-Validation}

\begin{frame}{Choosing the degree: the wrong way}
  \begin{columns}[T,onlytextwidth]
    \begin{column}{0.5\textwidth}
      \begin{block}{Tempting procedure}
        \footnotesize
        Fit $d=1,2,\dots,10$ on the training set. Compute $J_{\text{test}}$ for each. Pick
        the $d$ with the smallest $J_{\text{test}}$ and report that $J_{\text{test}}$.
      \end{block}
      \begin{pitfall}[The test set has been used to choose]
        \footnotesize
        $d$ is now a parameter fitted to the test set, exactly as $\vw$ was fitted to the
        training set. The reported $J_{\text{test}}$ is the \emph{best of ten tries} and is
        optimistic. In our data it would pick $d=7$ ($0.0380$), the luckiest of several
        near-equal degrees.
      \end{pitfall}
    \end{column}
    \begin{column}{0.46\textwidth}
      \begin{defnbox}[Three-way split]
        \small
        \begin{itemize}
          \item \textbf{Training set} (e.g.\ 60\%): fit $\vw,b$.
          \item \textbf{Cross-validation set} (validation, dev set; 20\%): choose $d$,
                $\lambda$, the architecture.
          \item \textbf{Test set} (20\%): report once, at the end.
        \end{itemize}
      \end{defnbox}
      \footnotesize
      Each set has one job. Any decision taken by looking at a set makes that set's error an
      underestimate.
    \end{column}
  \end{columns}
\end{frame}

\begin{frame}{Choosing the degree: the right way}
  \begin{columns}[T,onlytextwidth]
    \begin{column}{0.52\textwidth}
      \centering
      \begin{tikzpicture}
        \begin{axis}[width=7.2cm,height=4.8cm,xmin=0.5,xmax=10.5,ymin=0,ymax=0.65,
          xlabel={degree $d$},ylabel={cost},xtick={1,...,10},
          tick label style={font=\scriptsize},label style={font=\scriptsize},
          legend style={font=\tiny,at={(0.03,0.97)},anchor=north west},legend cell align=left]
          \addplot[cublue,very thick,mark=*,mark size=1.4pt] coordinates
            {(1,0.0806)(2,0.0777)(3,0.0219)(4,0.0209)(5,0.0203)(6,0.0203)(7,0.0201)
             (8,0.0167)(9,0.0165)(10,0.0164)};
          \addlegendentry{$J_{\text{train}}$}
          \addplot[curose,very thick,mark=square*,mark size=1.4pt] coordinates
            {(1,0.1119)(2,0.0991)(3,0.0496)(4,0.0515)(5,0.0510)(6,0.0516)(7,0.0518)
             (8,0.1472)(9,0.3432)(10,0.6085)};
          \addlegendentry{$J_{\text{cv}}$}
          \draw[cuamber,dashed,thick] (axis cs:3,0) -- (axis cs:3,0.4);
          \node[font=\tiny,cuamber,fill=white,inner sep=1pt] at (axis cs:3.2,0.43) {$d=3$};
        \end{axis}
      \end{tikzpicture}
    \end{column}
    \begin{column}{0.44\textwidth}
      \footnotesize
      \begin{block}{Procedure}
        \begin{enumerate}
          \item Fit each degree on the training set.
          \item Compute $J_{\text{cv}}$ for each; pick the smallest: $d=3$,
                $J_{\text{cv}}=0.0496$.
          \item Report $J_{\text{test}}$ for $d=3$ only: $\mathbf{0.0429}$.
        \end{enumerate}
      \end{block}
      \begin{keyidea}[Two curves, two stories]
        $J_{\text{train}}$ only falls as $d$ grows: more flexibility always fits the
        training data better. $J_{\text{cv}}$ falls, flattens, then rises --- the classic
        U. The bottom of the U is the model to choose.
      \end{keyidea}
      Same procedure for anything you choose: the number of layers and units in a network,
      $\lambda$, the learning rate.
    \end{column}
  \end{columns}
\end{frame}

% ===============================================================
\section{Diagnosing Bias and Variance}

\begin{frame}{Bias and variance, read from two numbers}
  \footnotesize
  \renewcommand{\arraystretch}{1.3}
  \begin{tabular}{@{}L{2.4cm}L{3.0cm}L{3.2cm}L{5.0cm}@{}}
    \toprule
    Diagnosis & $J_{\text{train}}$ & $J_{\text{cv}}$ & Our example\\
    \midrule
    \textbf{High bias} (underfit) & high & high, $\approx J_{\text{train}}$
      & $d=1$: $0.0806$ and $0.1119$\\
    \textbf{High variance} (overfit) & low & $\gg J_{\text{train}}$
      & $d=10$: $0.0164$ and $0.6085$\\
    \textbf{Just right} & low & low, near $J_{\text{train}}$
      & $d=3$: $0.0219$ and $0.0496$\\
    \textbf{Both} & high & $\gg J_{\text{train}}$
      & rarer with one feature; common in networks\\
    \bottomrule
  \end{tabular}
  \vspace{3mm}
  \begin{columns}[T,onlytextwidth]
    \begin{column}{0.5\textwidth}
      \begin{defnbox}[Bias]
        \small The model is too simple to represent the pattern, so it does poorly even on
        the data it was trained on. More of the same data will not help.
      \end{defnbox}
    \end{column}
    \begin{column}{0.46\textwidth}
      \begin{defnbox}[Variance]
        \small The model is so flexible that it fits the noise in its particular training
        set. A different training set would give a very different model.
      \end{defnbox}
    \end{column}
  \end{columns}
  \vspace{1mm}
  {\color{cugray} ``High'' and ``low'' need a reference point: the baseline, later in this
  lecture.}
\end{frame}

\begin{frame}{Regularisation moves a model along the same axis}
  \begin{columns}[T,onlytextwidth]
    \begin{column}{0.52\textwidth}
      \centering {\scriptsize Degree 10, regularised: $J(\vw,b)+\frac{\lambda}{2m}\|\vw\|^2$}\par
      \begin{tikzpicture}
        \begin{axis}[width=7.2cm,height=4.7cm,xmode=log,xmin=5e-5,xmax=200,ymin=0,ymax=0.13,
          xlabel={$\lambda$},ylabel={cost},tick label style={font=\scriptsize},
          label style={font=\scriptsize},
          legend style={font=\tiny,at={(0.03,0.97)},anchor=north west},legend cell align=left]
          \addplot[cublue,very thick,mark=*,mark size=1.4pt] coordinates
            {(0.0001,0.0174)(0.001,0.0189)(0.01,0.0198)(0.1,0.0211)(1,0.0299)(10,0.0525)
             (100,0.0768)};
          \addlegendentry{$J_{\text{train}}$}
          \addplot[curose,very thick,mark=square*,mark size=1.4pt] coordinates
            {(0.0001,0.0749)(0.001,0.0555)(0.01,0.0538)(0.1,0.0560)(1,0.0687)(10,0.0965)
             (100,0.1171)};
          \addlegendentry{$J_{\text{cv}}$}
          \draw[cuamber,dashed,thick] (axis cs:0.01,0) -- (axis cs:0.01,0.085);
          \node[font=\tiny,cuamber,fill=white,inner sep=1pt] at (axis cs:0.01,0.092) {$\lambda=0.01$};
        \end{axis}
      \end{tikzpicture}
    \end{column}
    \begin{column}{0.44\textwidth}
      \footnotesize
      \begin{block}{Reading the curve}
        \textbf{Small $\lambda$} (left): the degree-10 polynomial is nearly free ---
        high variance. At $\lambda=0$, $J_{\text{cv}}=0.6085$, off the chart.\\[2pt]
        \textbf{Large $\lambda$} (right): weights are forced towards $0$, $f$ flattens
        --- high bias.\\[2pt]
        In between, $J_{\text{cv}}$ is smallest at $\lambda=0.01$ ($0.0538$); its test cost
        is $0.0397$.
      \end{block}
      \begin{keyidea}[Mirror image]
        Against $d$, variance grows to the right. Against $\lambda$, variance grows to the
        \emph{left}. Choose $\lambda$ exactly as $d$: try a range roughly doubling or
        tenfold, and take the smallest $J_{\text{cv}}$.
      \end{keyidea}
    \end{column}
  \end{columns}
\end{frame}

% ===============================================================
\section{Baselines and Learning Curves}

\begin{frame}{``High'' compared with what? Establish a baseline}
  \begin{columns}[T,onlytextwidth]
    \begin{column}{0.5\textwidth}
      \footnotesize
      \begin{block}{Speech recognition (illustrative)}
        Some audio is so noisy that people cannot transcribe it either.\par
        \vspace{1mm}
        \renewcommand{\arraystretch}{1.15}
        \begin{tabular}{@{}lccc@{}}
          \toprule
          & case A & case B & case C\\
          \midrule
          human error     & $10.6\%$ & $10.6\%$ & $10.6\%$\\
          training error  & $10.8\%$ & $15.0\%$ & $15.0\%$\\
          validation error & $14.8\%$ & $15.5\%$ & $19.7\%$\\
          \midrule
          bias gap       & $0.2$ & $4.4$ & $4.4$\\
          variance gap   & $4.0$ & $0.5$ & $4.7$\\
          \midrule
          diagnosis & variance & bias & both\\
          \bottomrule
        \end{tabular}
      \end{block}
    \end{column}
    \begin{column}{0.46\textwidth}
      \footnotesize
      \begin{defnbox}[Baseline level of performance]
        \small The error you can reasonably hope to reach: human-level performance, a
        competing system, or past experience.
      \end{defnbox}
      \begin{block}{Two gaps}
        \textbf{bias gap} $=$ training error $-$ baseline\\
        \textbf{variance gap} $=$ validation error $-$ training error
      \end{block}
      \begin{keyidea}[Our example has a perfect baseline]
        The noise has $\sigma=0.25$, so no model can beat
        $J=\sigma^2/2=0.03125$ on average. Degree 3's training cost, $0.0219$, is
        \emph{below} it: the fit has absorbed some noise.
      \end{keyidea}
    \end{column}
  \end{columns}
\end{frame}

\begin{frame}{Learning curves: error against training-set size}
  \centering
  {\scriptsize Median over 100 random training sets of each size; validation set of 200}\par
  \begin{tikzpicture}
    \begin{axis}[name=lb,width=6.6cm,height=4.4cm,xmode=log,xmin=8,xmax=500,ymin=0,ymax=0.2,
      title={\footnotesize degree 1: high bias},xlabel={$m_{\text{train}}$},ylabel={cost},
      tick label style={font=\scriptsize},label style={font=\scriptsize},
      xtick={10,20,50,100,200,400},xticklabels={10,20,50,100,200,400},
      scaled y ticks=false,ytick={0,0.05,0.1,0.15,0.2},
      yticklabel style={/pgf/number format/fixed},
      legend style={font=\tiny,at={(0.97,0.5)},anchor=north east},legend cell align=left]
      \addplot[cublue,very thick,mark=*,mark size=1.2pt] coordinates
        {(10,0.0805)(12,0.0954)(15,0.1001)(20,0.1033)(30,0.1156)(50,0.1111)(100,0.1141)
         (200,0.1176)(400,0.1186)};
      \addlegendentry{$J_{\text{train}}$}
      \addplot[curose,very thick,mark=square*,mark size=1.2pt] coordinates
        {(10,0.1331)(12,0.1362)(15,0.1268)(20,0.1234)(30,0.1201)(50,0.1179)(100,0.1169)
         (200,0.1162)(400,0.1156)};
      \addlegendentry{$J_{\text{cv}}$}
      \addplot[cugray,dashed,domain=8:500,samples=2] {0.03125};
      \addlegendentry{noise floor}
    \end{axis}
    \begin{axis}[at={(lb.east)},anchor=west,xshift=10mm,width=6.6cm,height=4.4cm,
      xmode=log,xmin=15,xmax=500,ymin=0,ymax=0.2,
      title={\footnotesize degree 8: high variance},xlabel={$m_{\text{train}}$},
      xtick={20,50,100,200,400},xticklabels={20,50,100,200,400},
      scaled y ticks=false,ytick={0,0.05,0.1,0.15,0.2},
      yticklabel style={/pgf/number format/fixed},
      tick label style={font=\scriptsize},label style={font=\scriptsize}]
      \addplot[cublue,very thick,mark=*,mark size=1.2pt] coordinates
        {(20,0.0151)(30,0.0205)(50,0.0234)(100,0.0280)(200,0.0305)(400,0.0311)};
      \addplot[curose,very thick,mark=square*,mark size=1.2pt] coordinates
        {(20,0.1831)(30,0.0645)(50,0.0470)(100,0.0402)(200,0.0374)(400,0.0365)};
      \addplot[cugray,dashed,domain=15:500,samples=2] {0.03125};
    \end{axis}
  \end{tikzpicture}\par
  \vspace{1mm}
  \footnotesize
  \begin{columns}[T,onlytextwidth]
    \begin{column}{0.5\textwidth}
      \textbf{High bias.} Both curves flatten at about $0.117$, far above the floor, and
      meet. From $m=100$ to $m=400$, $J_{\text{cv}}$ improves by $0.001$:
      \alert{more data does not help}.
    \end{column}
    \begin{column}{0.46\textwidth}
      \textbf{High variance.} A large gap at small $m$ that closes as $m$ grows;
      $J_{\text{cv}}$ falls from $0.183$ to $0.037$, approaching the floor:
      \alert{more data helps}.
    \end{column}
  \end{columns}
\end{frame}

% ===============================================================
\section{What to Try Next, and Neural Networks}

\begin{frame}{Deciding what to try next, revisited}
  \begin{columns}[T,onlytextwidth]
    \begin{column}{0.55\textwidth}
      \footnotesize
      \renewcommand{\arraystretch}{1.3}
      \begin{tabular}{@{}L{4.2cm}L{3.4cm}@{}}
        \toprule
        Action & Fixes\\
        \midrule
        Get more training examples & high variance\\
        Try a smaller set of features & high variance\\
        Increase $\lambda$ & high variance\\
        \midrule
        Get additional features & high bias\\
        Add polynomial features & high bias\\
        Decrease $\lambda$ & high bias\\
        \bottomrule
      \end{tabular}
    \end{column}
    \begin{column}{0.41\textwidth}\raggedright
      \begin{keyidea}[Diagnose first, then act]\raggedright
        Compute $J_{\text{train}}$, $J_{\text{cv}}$ and a baseline. Whichever gap is larger
        tells you which half of the table to use.
      \end{keyidea}
      \begin{pitfall}[Do not shrink the training set]
        \footnotesize\raggedright Fewer examples make $J_{\text{train}}$ look better (the left end of
        a learning curve) but never fix bias --- and they make variance worse.
      \end{pitfall}
    \end{column}
  \end{columns}
\end{frame}

\begin{frame}[fragile]{Bias, variance and neural networks}
  \begin{columns}[T,onlytextwidth]
    \begin{column}{0.52\textwidth}
      \centering
      \begin{tikzpicture}[
          q/.style={draw=cublue,fill=culight,rounded corners,align=center,font=\scriptsize,
                    minimum width=3.3cm,minimum height=7mm},
          a/.style={draw=cuamber,fill=cuamber!12,rounded corners,align=center,font=\scriptsize,
                    minimum width=2.5cm,minimum height=7mm}]
        \node[q] (q1) at (0,0) {Does it do well on\\the training set?};
        \node[a] (a1) at (3.6,0) {bigger network};
        \node[q] (q2) at (0,-1.6) {Does it do well on\\the validation set?};
        \node[a] (a2) at (3.6,-1.6) {more data};
        \node[draw=cuteal,fill=cuteal!12,rounded corners,font=\scriptsize] (done) at (0,-3.0) {done};
        \draw[synapse,thick] (q1) -- node[above,font=\tiny]{no} (a1);
        \draw[synapse,thick] (a1.north) to[bend right=40] (q1.north);
        \draw[synapse,thick] (q1) -- node[left,font=\tiny]{yes} (q2);
        \draw[synapse,thick] (q2) -- node[above,font=\tiny]{no} (a2);
        \draw[synapse,thick] (a2.south) |- (-2.2,-3.55) |- (q1.west);
        \draw[synapse,thick] (q2) -- node[left,font=\tiny]{yes} (done);
      \end{tikzpicture}
    \end{column}
    \begin{column}{0.44\textwidth}
      \footnotesize
      \begin{block}{A recipe that rarely existed before}
        Large networks trained on modest data are usually \emph{low-bias}. So: grow the
        network until training error is acceptable, then add data until validation error
        is. Bias and variance can be attacked one at a time.
      \end{block}
      \begin{keyidea}
        A larger, well-regularised network almost never does worse than a smaller one. Its
        costs are computation and data, not accuracy.
      \end{keyidea}
\begin{lstlisting}
Dense(units=25, activation='relu',
      kernel_regularizer=
          tf.keras.regularizers.L2(0.01))
\end{lstlisting}
    \end{column}
  \end{columns}
\end{frame}

% ===============================================================
\section{Wrap-up and Assessment}

\begin{frame}{Summary}
  \begin{columns}[T,onlytextwidth]
    \begin{column}{0.5\textwidth}
      \begin{block}{What to remember}
        \footnotesize
        \begin{itemize}
          \item Judge a model on data it was not trained on; $J_{\text{train}}$ only
                measures fit.
          \item Train / cross-validation / test: fit, choose, report. Never choose on the
                test set.
          \item High bias: $J_{\text{train}}$ high. High variance:
                $J_{\text{cv}}\gg J_{\text{train}}$. Measure both against a baseline.
          \item $J_{\text{cv}}$ is U-shaped in $d$ and in $\lambda$ (mirrored).
          \item Learning curves: bias plateaus, more data does not help; variance gap
                closes, more data helps.
          \item Neural networks: bigger network for bias, more data for variance.
        \end{itemize}
      \end{block}
    \end{column}
    \begin{column}{0.46\textwidth}
      \begin{block}{Formulas}
        \[ J_{\text{cv}}=\frac{1}{2m_{\text{cv}}}\sum_{i}\bigl(f(\vx_{\text{cv}}^{(i)})-y_{\text{cv}}^{(i)}\bigr)^2 \]
        \[ \text{bias gap}=J_{\text{train}}-\text{baseline} \]
        \[ \text{variance gap}=J_{\text{cv}}-J_{\text{train}} \]
      \end{block}
      \begin{block}{Next}
        \footnotesize Part 2: the iterative loop of development, error analysis, adding
        data, transfer learning, deployment, fairness, and precision and recall.
      \end{block}
    \end{column}
  \end{columns}
\end{frame}

\begin{frame}{Descriptive questions}
  \footnotesize
  \begin{enumerate}
    \item List six actions you might take when a model makes large errors. Explain why a
          diagnostic should come first. \hfill\clo{5}
    \item Why is $J_{\text{train}}$ a poor estimate of a model's performance on new data?
          Define $J_{\text{test}}$ for regression and the misclassification error for
          classification. \hfill\clo{5}
    \item Explain why selecting the polynomial degree with the test set makes the reported
          test error optimistic. State the role of each part of a three-way split.
          \hfill\clo{6}
    \item Sketch $J_{\text{train}}$ and $J_{\text{cv}}$ against polynomial degree, label
          the high-bias and high-variance regions, and explain the shape of each curve.
          \hfill\clo{6}
    \item Using the slide's numbers, diagnose degrees 1, 3 and 10 of the running example.
          \hfill\clo{5}
    \item Sketch $J_{\text{train}}$ and $J_{\text{cv}}$ against $\lambda$ and explain why
          it is a mirror image of the degree plot. \hfill\clo{6}
  \end{enumerate}
\end{frame}

\begin{frame}{Descriptive questions (continued)}
  \footnotesize
  \begin{enumerate}\setcounter{enumi}{6}
    \item Define a baseline level of performance and give three ways of obtaining one.
          \hfill\clo{4}
    \item Human error is $2.0\%$, training error $2.3\%$, validation error $7.1\%$. Compute
          both gaps, give the diagnosis, and name two remedies. \hfill\clo{6}
    \item Draw typical learning curves for a high-bias and a high-variance model. For each,
          state whether collecting more data will help, and why. \hfill\clo{8}
    \item For each action in the ``what to try next'' table, state which problem it
          addresses and explain why. \hfill\clo{6}
    \item Describe the bias/variance recipe for neural networks and explain why it did not
          apply to small, classical models. \hfill\clo{5}
    \item In the running example, degree 3 reaches $J_{\text{train}}=0.0219$, below the
          noise floor of $0.03125$. Explain how this is possible and what it implies.
          \hfill\clo{6}
  \end{enumerate}
\end{frame}

\begin{frame}{Design questions}
  \footnotesize
  \begin{enumerate}
    \item \textbf{Design an evaluation protocol.} A Chattogram hospital has $3000$ labelled
          chest X-rays from $900$ patients and wants to choose among five network
          architectures and three values of $\lambda$. Design the split (including how you
          keep one patient's images together), the selection procedure, and what you report.
          \hfill\clo{12}
    \item \textbf{Design a diagnosis.} A crop-yield model has training $R^2=0.62$ and
          validation $R^2=0.60$, while an agronomist's estimate achieves $0.75$. Diagnose
          the model, and design the next three experiments in priority order, with the
          result that would make you stop each one. \hfill\clo{10}
    \item \textbf{Design a data-collection decision.} Collecting $1000$ more labelled
          examples will cost Tk~$5$ lakh. Design an experiment, using only the data you
          already have, that predicts whether it is worth paying for. \hfill\clo{12}
    \item \textbf{Design a $\lambda$ search.} Design a procedure to choose $\lambda$ for a
          network whose training takes four hours, under a one-week compute budget. Justify
          your grid and how you will use partial results. \hfill\clo{10}
  \end{enumerate}
\end{frame}

\begin{frame}{Reading}
  \begin{block}{Primary text}
    G.~James, D.~Witten, T.~Hastie and R.~Tibshirani,
    \emph{An Introduction to Statistical Learning}, 2nd ed., Springer, 2021.\\[2pt]
    \hfill{\small$\rightarrow$ \textbf{\S2.2.2} --- the bias--variance trade-off;
    \textbf{\S5.1} --- the validation-set approach and $k$-fold cross-validation;
    \textbf{\S6.2.1} --- choosing the ridge $\lambda$.}
  \end{block}
  \begin{block}{Supplementary}
    C.~M.~Bishop, \emph{Pattern Recognition and Machine Learning}, Springer, 2006.\\[2pt]
    \hfill{\small$\rightarrow$ \textbf{\S1.3} --- model selection and cross-validation;
    \textbf{\S3.2} --- the bias--variance decomposition.}
  \end{block}
  \vspace{1mm}
  \begin{center}
    \small\color{cugray}
    Lab (CSE 816, Week 6): model evaluation and selection; diagnosing bias and variance.
  \end{center}
\end{frame}

\end{document}
